\chapter{INTRODUCCI\'ON}

  La explosión de la información, así como la aparición de una gran variedad de
  formatos en los cuales se presenta, son dos características importantes de la
  llamada isociedad de la información, es en este contexto en el cual resulta
  necesario que las personas, las cuales requieren de información a diario,
  desarrollen habilidades de acceso, evaluación, uso y comunicación de la
  información de tal modo que sean capaces de aprovecharla y utilizarla
  eficazmente en diversos aspectos de su vidas.

  Ante esta situación, las bibliotecas tienen un rol que cumplir pues el
  carácter educativo que las identifica puede ser beneficioso si los
  bibliotecarios trabajan de manera participativa en la formación de los
  estudiantes colaborando así con el éxito académico. El rol que debe asumir el
  bibliotecario se centra en lograr que los estudiantes posean competencias
  informativas teniendo presente que un alumno que posee habilidades de
  búsqueda, análisis y uso de la información estará en ventaja sobre otros, lo
  cual le permitirá, sin duda, desarrollarse con éxito.

  Ante esta situación, las bibliotecas tienen un rol que cumplir pues el
  carácter educativo que las identifica puede ser beneficioso si los
  bibliotecarios trabajan de manera participativa en la formación de los
  estudiantes colaborando así con el éxito académico. El rol que debe asumir el
  bibliotecario se centra en lograr que los estudiantes posean competencias
  informativas teniendo presente que un alumno que posee habilidades de
  búsqueda, análisis y uso de la información estará en ventaja sobre otros, lo
  cual le permitirá, sin duda, desarrollarse con éxito.

  Ante esta situación, las bibliotecas tienen un rol que cumplir pues el
  carácter educativo que las identifica puede ser beneficioso si los
  bibliotecarios trabajan de manera participativa en la formación de los
  estudiantes colaborando así con el éxito académico. El rol que debe asumir el
  bibliotecario se centra en lograr que los estudiantes posean competencias
  informativas teniendo presente que un alumno que posee habilidades de
  búsqueda, análisis y uso de la información estará en ventaja sobre otros, lo
  cual le permitirá, sin duda, desarrollarse con éxito.

  Continúe aquí el texto xxxxx
  
Inicie aquí el texto, utilizando sangría de 1.25 cm. en el primer
párrafo. Continúe el segundo párrafo.

Continúe el segundo párrafo

\section{Primer subtítulo}

En la actualidad, las unidades de información hacen frente a muchos cambios
debido a los avances tecnológicos, explosión informativa, nuevos recursos y
soportes, por lo cual se implementan servicios innovadores que les permitan a
sus usuarios tener acceso a muchas fuentes de información. Para asegurar el
acceso y uso de los servicios los usuarios requieren poseer una serie de
habilidades que les permitan identificar, recuperar, manejar, discernir,
organizar, utilizar y comunicar la información de manera eficaz para la toma de
decisiones \cite{Reumann2012}.

Entre las causas se pueden considerar la falta de tiempo asignado al taller,
las limitaciones en cuanto a laboratorios, el poco personal, la falta de una
correcta selección de los contenidos, así como, una calificación que asegure el
cumplimiento de los objetivos.


\section{Segundo subtítulo}

\subsection{División del segundo subtítulo}

En la actualidad, las unidades de información hacen frente a muchos cambios
debido a los avances tecnológicos, explosión informativa, nuevos recursos y
soportes, por lo cual se implementan servicios innovadores que les permitan a
sus usuarios tener acceso a muchas fuentes de información. Para asegurar el
acceso y uso de los servicios los usuarios requieren poseer una serie de
habilidades que les permitan identificar, recuperar, manejar, discernir,
organizar, utilizar y comunicar la información de manera eficaz para la toma de
decisiones.

\begin{equation}
  \label{eq:1}
  a+b=\sqrt{\frac{4}{3}}
\end{equation}

Entre las causas se pueden considerar la falta de tiempo asignado al taller,
las limitaciones en cuanto a laboratorios, el poco personal, la falta de una
correcta selección de los contenidos, así como, una calificación que asegure el
cumplimiento de los objetivos.

\begin{figure}%[h]
  \centering
%  \includegraphics{images/testfig.png}
%  \includegraphics{images/testfig.png}
  \includegraphics[width=0.45\textwidth]{images/ELECTRONICA}
  \includegraphics[width=0.45\textwidth]{images/ELECTRONICA}
  \captionsetup{font=footnotesize}
  \caption{Scheme showing the architecture of a generic kinematic task.}
  \label{fig:diagram2}
\end{figure}



\section{Tercer subtítulo}

\subsection{División del tercer subtítulo}

En la actualidad, las unidades de información hacen frente a muchos cambios
debido a los avances tecnológicos, explosión informativa, nuevos recursos y
soportes, por lo cual se implementan servicios innovadores que les permitan a
sus usuarios tener acceso a muchas fuentes de información. Para asegurar el
acceso y uso de los servicios los usuarios requieren poseer una serie de
habilidades que les permitan identificar, recuperar, manejar, discernir,
organizar, utilizar y comunicar la información de manera eficaz para la toma de
decisiones.

Entre las causas se pueden considerar la falta de tiempo asignado al taller,
las limitaciones en cuanto a laboratorios, el poco personal, la falta de una
correcta selección de los contenidos, así como, una calificación que asegure el
cumplimiento de los objetivos.

